% !TeX program = tectonic
\documentclass[11pt,a4paper]{article}
\usepackage{fontspec}
\usepackage[english]{babel}
\babelfont{rm}[Language=Default]{Liberation Serif}
\babelfont[Language=Default]{sf}{Carlito}
\babelfont[Language=Default]{tt}{DejaVu Sans Mono}
\usepackage{amsmath,amssymb,amsthm}
\usepackage{geometry}
\geometry{a4paper, top=2.4cm, bottom=2.4cm, left=2.4cm, right=2.4cm}
\usepackage{booktabs,tabularx,enumitem,xcolor,tcolorbox}
\tcbuselibrary{breakable,skins}
\definecolor{linkblue}{HTML}{1F4E79}
\definecolor{statgreen}{HTML}{2F6B3A}
\definecolor{goldacc}{HTML}{8B7E5A}
\usepackage[colorlinks=true,linkcolor=linkblue,urlcolor=linkblue,unicode=true]{hyperref}
\theoremstyle{plain}
\newtheorem{theorem}{Theorem}
\newtheorem{lemma}[theorem]{Lemma}
\newtheorem{proposition}[theorem]{Proposition}
\newtheorem{corollary}[theorem]{Corollary}
\theoremstyle{definition}
\newtheorem{definition}[theorem]{Definition}
\theoremstyle{remark}
\newtheorem{remark}[theorem]{Remark}
\newtcolorbox{statusbox}[1][]{colback=statgreen!5,colframe=statgreen!75,
  title={\textbf{Summary of results:} #1},fonttitle=\small,boxrule=0.7pt,arc=2pt,breakable}
\newtcolorbox{databox}[1][]{colback=goldacc!6,colframe=goldacc!80,
  title={\textbf{Protocol data:} #1},fonttitle=\small,boxrule=0.7pt,arc=2pt,breakable}
\newtcolorbox{verifybox}[1][]{colback=linkblue!5,colframe=linkblue!80,
  title={\textbf{Verification:} #1},fonttitle=\small,boxrule=0.7pt,arc=2pt,breakable}
\widowpenalty=10000
\clubpenalty=10000
\setlength{\parskip}{0.35em}
\newcommand{\CC}{\mathbb{C}}
\newcommand{\RR}{\mathbb{R}}
\newcommand{\ZZ}{\mathbb{Z}}
\newcommand{\QQ}{\mathbb{Q}}
\newcommand{\Ga}{\Gamma}
\newcommand{\Bfun}{\mathrm{B}}
\newcommand{\im}{\operatorname{Im}}
\newcommand{\re}{\operatorname{Re}}
\newcommand{\lcm}{\operatorname{lcm}}
\newcommand{\SNF}{\operatorname{SNF}}
\newcommand{\Jac}{\operatorname{Jac}}
\newcommand{\rank}{\operatorname{rank}}
\newcommand{\Ind}{\operatorname{Index}}
\newcommand{\disc}{\operatorname{disc}}
\begin{document}
\begin{center}
{\LARGE\bfseries The K3 Stand: Neron–Severi Rank 20}\\[0.4em]
{\large forty-eight lines and the exact intersection matrix}\\[0.6em]
{\normalsize Isaev Iskhak Khamzatovich}\\[0.2em]
{\normalsize The <<Dynamic Principle>> Program --- the \texttt{hodge-laboratory}}\\[0.15em]
{\small Standalone theorem monograph · Монография 8/16}
\end{center}
\vspace{0.4em}
{\color{linkblue}\hrule height 0.8pt}
\vspace{0.8em}

\section{Setting}

The Fermat quartic $x^4+y^4+z^4+w^4=0\subset\mathbb{P}^3$ is a K3
surface of maximal N\'eron--Severi rank. The Hodge conjecture for it
is proved (Shioda, 1979); the program reproduces all invariants by
exact integer arithmetic over $\QQ$.

\begin{theorem}[invariants of the K3 stand]\label{th:k3stand}
For the exact intersection matrix $49\times49$ (48 lines + the
hyperplane class $h$): (1) $\rank_\QQ G=20$ --- matches Shioda's
$\rho=20$; (2) the signature is $(1,19)$; (3) the sublattice
determinant is $64=8^2$; (4) the sum of the four lines of one family
is equivalent to $h$.
\end{theorem}

\begin{proof}
The lines: three families of $16$: $L_1(a,b)\colon x+i^ay=0$,
$z+i^bw=0$; $L_2$: $x+i^az=0$, $y+i^bw=0$; $L_3$: $x+i^aw=0$,
$y+i^bz=0$, $a,b\in\ZZ/4$. The rules are combinatorial:
$L\cdot L=-2$; within a family $1$ iff exactly one index matches;
between families cyclic conditions ($L_1\cdot L_2$:
$a_1+b_2\equiv a_2+b_1$; $L_2\cdot L_3$: $a_1+b_1\equiv a_2+b_2$;
$L_1\cdot L_3$: $c\equiv a+b+d+1$ --- the extra unit from three odd
exponents); $h\cdot L=1$, $h^2=4$. Rank: exact fraction-free
elimination --- $20$ nonzero rows; double registration via Smith
minors. Signature: one positive, $19$ negative eigenvalues, none
zero (separation threshold $0.1$). Determinant: the full K3 lattice
is unimodular of signature $(3,19)$; the sublattice index is $8$,
$\det=64=8^2$. Equivalence: the intersection vector of the sum
equals that of $h$ componentwise (all $49$); the difference is
orthogonal to the algebraic sublattice and lies in
$H^{2,0}\oplus H^{0,2}$ by the Hodge index theorem; but it is
algebraic, and algebraic classes pair trivially with $H^{2,0}$;
hence zero.
\end{proof}

\begin{proposition}[codimension-2 cycles]
The point class: $[\mathrm{pt}]=h^2/4$ --- exactly, since
$(h^2)^2=\deg X=4$ and $H^4(X)=\QQ h^2$. On $K3\times K3$: the pairs
$(c_2,h_2)$ give $24$ independent codim-$(2,2)$ classes, the
diagonal $h^2$-pairs give $4$ more; all algebraic with explicit
supports. No blind spot: every $(2,2)$ class orthogonal to them is
zero by the Hodge index theorem.
\end{proposition}

\begin{databox}[numbers]
$n_{\mathrm{lines}}=48$; $\rank=20=\rho$; inertia $(1,19)$;
$\det=64$; hyperplane relation exact; $12$ orbits; character
multiplicities $(1,7,7,7)$ summing to $22$; codim 2: $4+24$;
runtime $1.04$~s.
\end{databox}

\begin{statusbox}[summary]
All K3 invariants are computed exactly over $\QQ$; the agreement
with Shioda is confirmed; the family--hyperplane equivalence is
proved exactly.
\end{statusbox}

\begin{verifybox}[reproduction]
\texttt{python3 laboratory.py} $\to$ ``Stands $\to$ K3'': the
matrix, rank, signature, determinant --- within a second.
\end{verifybox}

\vfill
{\color{linkblue}\hrule height 0.6pt}
\vspace{0.5em}
{\small\color{gray}
License: individual exclusive license of Isaev Iskhak Khamzatovich.
All rights to the computations, formulas, and texts belong to the author.
Verification: \texttt{python3 laboratory.py} (``Stands''/``Certificates''),
Lean: \texttt{verification/lean/HodgeLaboratory.lean}.
}
\end{document}
